\documentclass{beamer}

\usepackage[utf8]{inputenc}
\input{../helper}

\title{Introdução à Computação \\ redimencionando arrays}
\author{Alexandre Rademaker}
\date{}

\begin{document}

\begin{frame}
 \maketitle
\end{frame}

\begin{frame}[fragile,allowframebreaks]{arrays}

  Um array é uma porção contínua na memória ocupada por segmentos de
  igual tamanho, logo mesmo tipo.

  Quando criamos um array, definimos sua capacidade de armazenamento,
  isto é, seu tamanho.

  {\color{red} Lembrem-se} que em C tentar acessar uma posição além do
  limite do array não gera erro!

  Mas e se no decorrer da computação precisamos de mais espaço?

  \framebreak

  Temos duas possibilidades.

  Alocar um espaço maior do que o necessário de início, eventualmente
  desperdiçando espaço. Mas precisamos definor o pior caso, o máximo
  que ``podemos precisar''.

  Quando o limite for atingido, criar um novo e maior array, copiar os
  valores para o novo array.

\end{frame}


\begin{frame}[fragile,allowframebreaks]{redimencionando arrays}

  Usando alocação estática 

  \begin{lstlisting}[style=normalc,caption=src/list0.c]
    int list[3];
    list[0] = 1; list[1] = 2; list[2] = 3;
    
    // ops! I need 10 slots
    
    int new[10];
    for (int i = 0; i < 10; i++) 
      new[i] = i < 3 ? list[i] : 0;
  \end{lstlisting}

  \framebreak

  Se quisermos alocar o array no \textbf{heap} 

  \begin{lstlisting}[style=normalc,caption=src/list1.c]
    int *list = malloc(3 * sizeof(int));
    if (list == NULL) return 1;
    
    list[0] = 1; list[1] = 2; list[2] = 3;

    // ops! I need 10 slots
    
    int *new = malloc(10 * sizeof(int));
    if (new == NULL) { free(list); return 1; }
    
    for (int i = 0; i < 10; i++) 
      new[i] = i < 3 ? list[i] : 0;

    free(list); // important
    list = new;  
  \end{lstlisting}

  \framebreak

  Usando alocação estática, não podemos renomear um array

  {\color{red} \texttt{list = new;}}

  para isso precisamos de ponteiros. Em \texttt{list0.c} vejamos como
  a memória é ocupada.

  Após transferir valores, com ponteiros, precisamos lembrar de
  liberar a memória antes de 'atualizar' o ponteiro.

\end{frame}


\begin{frame}[fragile,allowframebreaks]{redimencionando arrays}

  Mas o processo pode ser simplificado com
  \href{https://linux.die.net/man/3/realloc}{realloc} da biblioteca
  \texttt{stdlib.h}, vide \texttt{man 3 realloc}.

  \texttt{void *realloc(void *ptr, size\_t size);}

  O tipo \texttt{size\_t} é definido como {\color{teal} \texttt{unsigned integer}}.

  Mas atenção, as novas posições podem não ser `zeradas'.

  \framebreak

  \begin{lstlisting}[style=normalc,caption=src/list2.c]
    int *list = malloc(3 * sizeof(int));
    if (list == NULL)
      return 1;

    list[0] = 1; list[1] = 2; list[2] = 3;

    list = realloc(list, 10 * sizeof(int));
    if (list == NULL) {
     free(list); return 1;
    }

    for (int i = 3; i < 10; i++)
      list[i] = 0;

    free(list); return 0;
  \end{lstlisting}
  
\end{frame}


\begin{frame}[allowframebreaks]{operações com arrays}

  Para recuperar o valor de índice $n$ (acesso aleatório) no array, a
  complexidade é $O(1)$.

  Se o array estiver ordenado, a complexidade de localizar se um valor
  existe no array é $O(\log n)$ (digamos que queremos adicionar um
  elemento que não exista).

  E temos um tempo $O(n)$ para adicionar um novo valor (copiando um
  por um para uma nova região) se o array estiver 'lotado'.

  O melhor caso, para inserção e pesquisa têm $\Omega(1)$, se
  encontrarmos o valor imediatamente na primeira posição, ou ter
  memória livre após nosso array para adicionar um novo valor.

\end{frame}


\end{document}

%%% Local Variables:
%%% mode: latex
%%% TeX-master: t
%%% End:
